\documentclass[10pt,a4paper]{amsart}
\usepackage[margin=19mm]{geometry}
\usepackage{amsmath,amssymb,mathtools}
\usepackage[scr=rsfs,cal=euler]{mathalfa}
\usepackage{xfrac}
\usepackage[unicode,hidelinks,bookmarks=false]{hyperref}
\newcommand{\ie}{i.e.\ }
\newcommand{\cf}{cf.\ }
\newcommand{\resp}{respectively\ }
\newcommand{\Subsection}{\subsection}
\providecommand{\nobreakdash}{\nobreak\hbox{-}}
\setlength{\emergencystretch}{2em}
\multlinegap0pt
\newtheorem{thm}{Theorem}
\newtheorem{cor}[thm]{Corollary}
\newtheorem{lem}[thm]{Lemma}
\newtheorem{prop}[thm]{Proposition}
\title[Theorem 4.11]{Morse index and nullity of rotationally symmetric $Y$-singular minimal surfaces: the resonant surface $Y_{\pi/6}$ (Theorem 4.11)}
\author{Elham Matinpour}
\date{Source: arXiv:2410.15534v2 (CC BY 4.0)}
\begin{document}
\maketitle

\noindent\emph{Source and licence.} Morse index and nullity of rotationally symmetric $Y$-singular minimal surfaces, by Elham Matinpour. Version arXiv:2410.15534v2
(CC BY 4.0), distributed by arXiv under the Creative Commons Attribution 4.0 International licence,
http://creativecommons.org/licenses/by/4.0/, as recorded in the arXiv licence metadata for this exact
version and shown on the abstract page https://arxiv.org/abs/2410.15534v2. Full licence notice:
\texttt{licenses/arxiv-2410.15534-CC-BY-4.0.txt}.\par\smallskip

\noindent\textit{Editorial scope.} Counted unit: the article's Theorem 4.11 (source0.tex lines 1728--1730, printed on page 30 of the published version): the surface $Y_{\pi/6}$ has Morse index two and nullity five. The article's complete original proof of it is delivered with it (source0.tex lines 1732--1930, one \texttt{proof} environment of 199 lines, adjacent to the statement, with no gap and no deferral). No mathematics is written, repaired or re-derived: the statement and the proof are the article's own lines, in the article's own notation, no retained line is substituted, re-flowed or omitted, and the package carries no omission record because nothing was removed. Retained uncounted as the source-local context the proof needs, each exactly the statement of the result it invokes and nothing else: the section and subsection headings that introduce the framework (lines 192--195, printed in the article as Section 3 and Subsection 3.1); the closedness and orthogonality result for the admissible spaces (lines 646--675, printed as Proposition 3.5, source label \texttt{prop: second assumption steklov}); the general index decomposition at a nonresonant frequency (lines 789--827, printed as Corollary 3.6, source label \texttt{coro: general index decom}); the noncompact variational framework (lines 891--943, printed as Proposition 3.7, source label \texttt{prop: function space for noncompact index}); the kernel of the Jacobi operator of catenoidal faces (lines 1094--1129, printed as Proposition 4.1, source label \texttt{prop: kernel(J) of cat-faces}); the zero-junction-trace Dirichlet Morse index of a catenoidal end (lines 1451--1461, printed as Lemma 4.7, source label \texttt{lem: dirichlet index catenoidal faces}); and the heading of the subsection in which the counted theorem stands (line 1727, printed in the article as Subsection 4.6). Every cross-reference of every retained line resolves inside this document, and no retained line contains a citation, so the wrapper carries no bibliography and no bibliography entry. No retained line contains a figure environment, an image-inclusion command or drawing code, so no image is delivered and the package declares no asset dependency; the single figure of the article is not retained.

Editorial formatting only, in addition: an AMS \texttt{amsart} preamble that restates the article's own declarations used by the retained lines -- the four theorem-like environments \texttt{thm}, \texttt{cor}, \texttt{lem} and \texttt{prop}, kept on one shared counter exactly as in the article but without the article's section-based prefix -- and the article's own \texttt{amsart} class, which is the class the article itself uses. Every display of the retained lines is an unnumbered \texttt{\textbackslash[}\texttt{\textbackslash]} display, so no equation number is renumbered. The retained section and subsection headings print in this document as Section 1, Subsection 1.1 and Subsection 1.2, whereas the article prints them as Section 3, Subsection 3.1 and Subsection 4.6; the retained environments print as Proposition 1, Corollary 2, Proposition 3, Proposition 4, Lemma 5 and Theorem 6 in order of appearance, whereas the article prints exactly the same six items as Proposition 3.5, Corollary 3.6, Proposition 3.7, Proposition 4.1, Lemma 4.7 and Theorem 4.11.

The complete native source of the article as distributed in the arXiv e-print for this exact version is bundled unchanged as \texttt{sources/167136/source0.tex}: it is byte identical to the e-print file \texttt{main.tex} except that the pipeline's mechanical concatenation prepends one header comment line naming it. No publisher class or style file is shipped in the e-print and none is redistributed or used.
\section{Variational Framework and Index Decomposition} \label{sec: index evaluation on catenoidal}
We now set up the variational framework used to compute the index of the $Y$-noids. The main points are the coupled boundary condition at the junction and the separation into angular modes.

\subsection{Index Decomposition} \label{subsec: index decomp}

\begin{prop}\label{prop: second assumption steklov}
For every $Y$-noid, the spaces $\mathcal D_0$, $\mathcal D_n^c$, and $\mathcal D_n^s$ are closed. Distinct angular sectors are orthogonal for the $\mathcal D(M)$ inner product and for $Q$. For every $u,v\in\mathcal D(M)$,
\[
S_Nu\longrightarrow u\quad\text{in }\mathcal D(M),
\]
\[
\|u\|_{\mathcal D(M)}^2
=\|P_0u\|_{\mathcal D(M)}^2
+\sum_{n=1}^\infty\left(\|P_n^cu\|_{\mathcal D(M)}^2+\|P_n^su\|_{\mathcal D(M)}^2\right),
\]
and
\[
Q(u,v)=Q(P_0u,P_0v)
+\sum_{n=1}^\infty\left(Q(P_n^cu,P_n^cv)+Q(P_n^su,P_n^sv)\right),
\]
with absolute convergence of the last series. Each projection preserves trace compatibility and the facewise Jacobi equation.

Set
\[
\iota_0=\operatorname{Ind}(Q;\mathcal D_0),
\qquad
\iota_n=\operatorname{Ind}(Q;\mathcal D_n^c)=\operatorname{Ind}(Q;\mathcal D_n^s),\quad n\ge1.
\]
Then
\[
\operatorname{Ind}(Q;\mathcal D(M))
=\iota_0+2\sum_{n=1}^\infty\iota_n.
\]
The equality is understood in $\{0,1,2,\ldots,\infty\}$. The weak kernel decomposes as the completed orthogonal sum of the restricted bilinear kernels in these angular sectors.
\end{prop}

\begin{cor}\label{coro: general index decom}
For every nonresonant frequency,
\[
\mathcal X_n=\mathcal X_n^0\oplus_{\mathfrak q_n}\mathscr E_n(V_\Delta),
\qquad
\mathfrak q_n(\mathscr E_n C,a)
=R_\Gamma\sum_{i=1}^3(\delta_{i,n}-\beta_i)C_i a_i^\Gamma.
\]
The decomposition is direct, and its projections are bounded in the radial variational norm. In particular,
\[
\iota_n=\sum_{i=1}^3d_{i,n}
+\operatorname{Ind}(\mathfrak b_n;V_\Delta).
\]
Its bilinear kernel is
\[
\ker\mathfrak q_n
=\mathscr E_n(\ker\mathfrak b_n),
\qquad
C\in\ker\mathfrak b_n
\ \Longleftrightarrow\
(\delta_{i,n}-\beta_i)C_i=r_n\quad(i=1,2,3)
\]
for some real $r_n$, with $C\in V_\Delta$. For the corresponding two-dimensional field, the common boundary function is $r(\theta)=r_n\phi(\theta)$, so it need not be constant along $\Gamma$. Here $\ker\mathfrak q_n$ is tested against all of $\mathcal X_n$, and $\ker\mathfrak b_n$ against all of $V_\Delta$.

For every $Y$-noid, including one whose radial sector is resonant,
\[
\operatorname{Ind}(Q;\mathcal D(M))
=\iota_0+2\sum_{n=1}^{\infty}
\left(\sum_{i=1}^3d_{i,n}+\operatorname{Ind}(\mathfrak b_n;V_\Delta)\right).
\]
If every catenoidal $T_i$ is nonzero, this also gives
\[
\operatorname{Ind}(Q;\mathcal D(M))
=\sum_{i=1}^3 \operatorname{Ind}_{J_i}^0(M_i)
+\operatorname{Ind}(\mathfrak b_0;V_\Delta)
+2\sum_{n=1}^{\infty}\operatorname{Ind}(\mathfrak b_n;V_\Delta).
\]
Here $\operatorname{Ind}_{J_i}^0(M_i)$ is the index on the zero-junction-trace subspace of $\mathcal D(M_i)$, equivalently on its compactly supported functions. Sums of indices are understood in $\{0,1,2,\ldots,\infty\}$.
\end{cor}

\begin{prop}[Noncompact variational framework]\label{prop: function space for noncompact index}
Let $M$ be one of the $Y$-noids considered in this paper, and let $\mathcal D(M)$ and $L_*^2(M)$ be as in Subsection~\ref{subsec: index decomp}. Then the following hold.
\begin{enumerate}
\item[(a)] $\mathcal D(M)$ is a Hilbert space, dense in $L_*^2(M)$. The index form $Q$ is continuous on $\mathcal D(M)$ and defines a closed, lower-semibounded form on $L_*^2(M)$.

\item[(b)] Compactly supported admissible $H^1$ functions are dense in $\mathcal D(M)$. Consequently,
\[
\operatorname{Ind}(Q;\mathcal D(M))
=
\operatorname{Ind}(Q;W_{\Delta,c}^{1,2}(M)).
\]
This equality concerns the Morse index only. The cutoff argument does not preserve the bilinear kernel.

\item[(c)] Let $\mathscr A_*$ be the self-adjoint operator associated with $Q$ on the weighted Hilbert space $L_*^2(M)$. Writing $f_i=u_i|_\Gamma$, its domain is
\[
\begin{split}
\operatorname{Dom}(\mathscr A_*)=\biggl\{u\in\mathcal D(M):\;&
\sum_{i=1}^3\int_{M_i}\omega^{-1}|J_i u_i|^2\,dA_i<\infty,\\
&\partial_{\nu_i}u_i-\beta_i f_i=r\quad(i=1,2,3)\\
&\text{for some }r\in H^{-1/2}(\Gamma)\biggr\}.
\end{split}
\]
On this domain,
\[
\mathscr A_*u=(-\omega^{-1}J_i u_i)_{i=1}^3.
\]
In particular,
\[
\ker\mathscr A_*=\ker_{\mathcal D}Q,
\]
and $u\in\ker_{\mathcal D}Q$ if and only if
\[
J_i u_i=0\quad\text{on }M_i,\qquad
\sum_{i=1}^3 f_i=0,\qquad
\partial_{\nu_i}u_i-\beta_i f_i=r\quad\text{on }\Gamma
\]
for one common $r\in H^{-1/2}(\Gamma)$.

\item[(d)] A nonzero-trace field $u\in\mathcal D(M)$ satisfies the weak coupled Steklov equation
\[
Q(u,v)=\sigma\sum_{i=1}^3\int_\Gamma f_i g_i\,ds
\qquad(v\in\mathcal D(M)),
\]
where $g_i=v_i|_\Gamma$, if and only if
\[
J_i u_i=0\quad\text{on }M_i,\qquad
\sum_{i=1}^3f_i=0,
\qquad
\partial_{\nu_i}u_i-\beta_i f_i-\sigma f_i=r\quad\text{on }\Gamma
\]
with the same $r\in H^{-1/2}(\Gamma)$ for all three faces.
\end{enumerate}
\end{prop}

\begin{prop} \label{prop: kernel(J) of cat-faces}
Let
\[
X(t,\theta)=c\big(\cosh(t+T)\cos\theta,\cosh(t+T)\sin\theta,\pm t\big),
\qquad t\in[0,\infty),
\]
parameterize a catenoidal end $\Sigma$, with $T\in\mathbb R$. Define
\[
\ker_{\mathcal D}J(\Sigma)=\{u\in\mathcal D(\Sigma):Ju=0\}.
\]
Then
\[
\ker_{\mathcal D}J(\Sigma)
=\overline{\bigoplus_{n\ge0}W_n}^{\,\mathcal D(\Sigma)},
\]
where
\[
\begin{aligned}
W_0&=\operatorname{span}\{\tanh(t+T)\},\\
W_1&=\operatorname{span}\left\{\frac{\cos\theta}{\cosh(t+T)},\frac{\sin\theta}{\cosh(t+T)}\right\},\\
W_n&=\operatorname{span}\left\{\cos(n\theta)(n+\tanh(t+T))e^{-n(t+T)},\right.\\
&\hspace{3.1cm}\left.\sin(n\theta)(n+\tanh(t+T))e^{-n(t+T)}\right\},\quad n\ge2.
\end{aligned}
\]
For $T\neq0$, the restrictions of these functions to $\{t=0\}$ are facewise Steklov eigenfunctions with
\[
\delta_0=-\frac1{c\cosh^2T\sinh T},\qquad
\delta_1=\frac{\sinh T}{c\cosh^2T},
\]
and
\[
\delta_n=-\frac{1-(n\cosh^2T)(n+\tanh T)}
{c(n+\tanh T)\cosh^3T},\qquad n\ge2.
\]
When $T=0$, the radial field $\tanh t$ has zero trace on the junction and hence is a Dirichlet Jacobi field rather than a nonzero Steklov trace.
\end{prop}

\begin{lem}\label{lem: dirichlet index catenoidal faces}
Let $\Sigma$ be a catenoidal end parameterized as in Proposition \ref{prop: kernel(J) of cat-faces}, with parameter $T$. Then its zero-junction-trace Dirichlet Morse index is
\[
\operatorname{Ind}_{J}^{0}(\Sigma)=
\begin{cases}
1,&T<0,\\
0,&T\ge0.
\end{cases}
\]
The zero-trace bilinear kernel is trivial when $T\ne0$ and is spanned by $\tanh t$ when $T=0$. Consequently, for the faces in the $Y$-noid family, every non-graphical catenoidal face has Dirichlet Morse index one and every graphical catenoidal face has Dirichlet Morse index zero.
\end{lem}

\subsection{The Resonant Surface \texorpdfstring{$Y_{\pi/6}$}{Y pi/6}}

\begin{thm}
The surface $Y_{\pi/6}$ has Morse index two and nullity five.
\end{thm}

\begin{proof}
Let $M=(\Sigma_1,\Sigma_2,\Sigma_3;\Gamma)=Y_{\pi/6}$. Here $\Sigma_1$ is non-graphical, $\Sigma_2,\Sigma_3$ are graphical, and
\[
\alpha_1=\frac{\pi}{6},\qquad
\alpha_2=\frac{5\pi}{6},\qquad
\alpha_3=-\frac{\pi}{2}.
\]
Hence $T_1,T_2\neq0$ and $T_3=0$. The radial sector is therefore resonant only on $\Sigma_3$.

\medskip
\noindent Step 1. \emph{Positive angular frequencies.}
For a zero-junction-trace profile in frequency $n=1$, the catenoidal one-dimensional form satisfies
\[
\mathfrak q_{i,1}(a_i,a_i)
=\int_0^\infty\left|a_i'+\tanh(t+T_i)a_i\right|^2\,dt\geq0.
\]
Indeed, expanding the square and integrating the cross term gives the $n=1$ form. The boundary term at $t=0$ vanishes because $a_i(0)=0$, and the term at infinity vanishes by the cutoff approximation in Proposition~\ref{prop: function space for noncompact index}. For $n\ge2$,
\[
n^2-2\operatorname{sech}^2(t+T_i)\ge n^2-2>0,
\]
so every positive-frequency zero-trace form is nonnegative.

For the boundary part, Corollary~\ref{coro: general index decom} gives
\[
\delta_{i,1}-\beta_i=0,
\qquad
\delta_{i,n}-\beta_i
=\frac{n^2-1}{R_\Gamma(n-\cos\alpha_i)}>0,
\quad n\ge2.
\]
Thus $\mathfrak b_1\equiv0$ and $\mathfrak b_n$ is positive definite for $n\ge2$. Consequently
\[
\iota_n=0\qquad(n\ge1).
\]
Hence all negative directions lie in the radial sector, and
\[
\operatorname{Ind}(Q;\mathcal D(M))=\iota_0.
\]

\medskip
\noindent Step 2. \emph{The zero-trace radial space.}
Let $\mathcal X_0$ be the radial variational space and $\mathcal X_0^0$ its zero-junction-trace subspace, as introduced before Corollary~\ref{coro: general index decom}. By Lemma~\ref{lem: dirichlet index catenoidal faces} the Dirichlet indices of $\Sigma_1,\Sigma_2,\Sigma_3$ are $1,0,0$. Step~1 shows that the positive-frequency zero-trace forms are nonnegative, so the unique zero-trace negative direction on $\Sigma_1$ is radial. Therefore
\[
\operatorname{Ind}(\mathfrak q_0;\mathcal X_0^0)=1.
\]

The bilinear kernel of $\mathfrak q_0$ on $\mathcal X_0^0$ is one-dimensional. Indeed, Proposition~\ref{prop: kernel(J) of cat-faces} shows that an admissible radial Jacobi field on a catenoidal face is a multiple of $\tanh(t+T_i)$. This profile has nonzero junction value when $T_i\neq0$ and has zero junction value when $T_i=0$. Hence only $\Sigma_3$ contributes a zero-trace radial Jacobi field, and
\[
\ker\bigl(\mathfrak q_0|_{\mathcal X_0^0}\bigr)
=\operatorname{span}\{k\},
\qquad
k=(0,0,\tanh t).
\]

For $i=1,2$, let
\[
p_i(t)=\frac{\tanh(t+T_i)}{\tanh T_i}
\]
be the admissible radial Jacobi profile with junction value one. Set
\[
e=(p_1,-p_2,0),
\qquad
\mathcal Y=\{a\in\mathcal X_0:a_3^\Gamma=0\}.
\]
The trace of every $a\in\mathcal Y$ is of the form $(s,-s,0)$. Thus, with $s=a_1^\Gamma$,
\[
a=(a-se)+se,
\qquad a-se\in\mathcal X_0^0,
\]
and the decomposition is unique. Hence
\[
\mathcal Y=\mathcal X_0^0\oplus\operatorname{span}\{e\}
\]
algebraically.

Because $p_1$ and $p_2$ are Jacobi profiles, the boundary-pairing identity gives
\[
\mathfrak q_0(e,w)=0
\qquad(w\in\mathcal X_0^0).
\]
Moreover,
\[
\delta_{i,0}-\beta_i=\frac{1}{R_\Gamma\cos\alpha_i},
\qquad i=1,2.
\]
Since $\cos(\pi/6)=-\cos(5\pi/6)$,
\[
(\delta_{1,0}-\beta_1)+(\delta_{2,0}-\beta_2)=0,
\]
and therefore $\mathfrak q_0(e,e)=0$. We obtain the orthogonal decomposition
\[
\mathcal Y
=\mathcal X_0^0\oplus_{\mathfrak q_0}\operatorname{span}\{e\},
\]
so
\[
\operatorname{Ind}(\mathfrak q_0;\mathcal Y)=1,
\qquad
\ker(\mathfrak q_0|_{\mathcal Y})=\operatorname{span}\{k,e\}.
\]

\medskip
\noindent Step 3. \emph{The remaining radial trace direction and the Morse index.}
The space $\mathcal Y$ has codimension one in $\mathcal X_0$. To add the remaining trace direction, take
\[
h(t)=\operatorname{sech}t
\]
on $\Sigma_3$ and define
\[
\psi=(0,-p_2,h).
\]
Since $h(0)=1$, the junction trace of $\psi$ is $(0,-1,1)$, and therefore
\[
\mathcal X_0=\mathcal Y\oplus\operatorname{span}\{\psi\}
\]
algebraically.

The key pairing with the resonant Dirichlet Jacobi field is finite. Since $k(0)=0$, the junction term vanishes, and only the $\Sigma_3$ interior form remains. Using $k(t)=\tanh t$ and the equation $k''+2\operatorname{sech}^2t\,k=0$,
\[
\begin{aligned}
\mathfrak q_0(k,\psi)
&=\int_0^\infty
\left(k'h'-2\operatorname{sech}^2t\,kh\right)\,dt\\
&=[k'h]_{0}^{\infty}=-1.
\end{aligned}
\]

Because $\mathcal Y$ has codimension one in $\mathcal X_0$, adjoining one vector can raise the index by at most one. More explicitly, if $L\subset\mathcal X_0$ is negative definite, then $\dim(L\cap\mathcal Y)\ge\dim L-1$. Hence
\[
\iota_0\le
\operatorname{Ind}(\mathfrak q_0;\mathcal Y)+1=2.
\]
For the reverse inequality, choose $v\in\mathcal Y$ with $\mathfrak q_0(v,v)<0$ and set
\[
\psi'
=\psi-
\frac{\mathfrak q_0(v,\psi)}{\mathfrak q_0(v,v)}v.
\]
Then
\[
\mathfrak q_0(v,\psi')=0,
\qquad
\mathfrak q_0(k,\psi')=-1,
\]
because $k$ lies in the bilinear kernel of $\mathfrak q_0|_{\mathcal Y}$. On $\operatorname{span}\{k,\psi'\}$ the matrix of the form is
\[
\begin{pmatrix}
0&-1\\
-1&\mathfrak q_0(\psi',\psi')
\end{pmatrix},
\]
whose determinant is $-1$. Hence this two-dimensional restriction has one positive and one negative direction. Its negative direction is $\mathfrak q_0$-orthogonal to the negative line $\operatorname{span}\{v\}$. These two negative directions span a negative-definite two-dimensional subspace of $\mathcal X_0$. Therefore
\[
\iota_0\ge2.
\]
Combining the two bounds gives
\[
\iota_0=2.
\]
Together with Step~1,
\[
\operatorname{Ind}(Y_{\pi/6})=2.
\]

\medskip
\noindent Step 4. \emph{The weak nullity.}
Let $u=y+s\psi\in\ker\mathfrak q_0$, with $y\in\mathcal Y$. Testing against $k$ gives
\[
0=\mathfrak q_0(u,k)
=s\,\mathfrak q_0(\psi,k)=-s,
\]
so $s=0$. Hence $u=y\in\ker(\mathfrak q_0|_{\mathcal Y})=\operatorname{span}\{k,e\}$. It remains to impose the equation against the one missing test direction $\psi$. The corresponding functional is nonzero because
\[
\mathfrak q_0(k,\psi)=-1.
\]
Therefore it cuts the two-dimensional space $\operatorname{span}\{k,e\}$ down to a one-dimensional radial weak kernel.

The radial kernel line can be written explicitly. The boundary-pairing identity gives
\[
\mathfrak q_0(e,\psi)
=R_\Gamma(\delta_{2,0}-\beta_2)
=\frac{1}{\cos\alpha_2}
=-\frac{2}{\sqrt3}.
\]
Thus
\[
\ker\mathfrak q_0
=\operatorname{span}\left\{e-\frac{2}{\sqrt3}k\right\}.
\]

For $n=1$, Corollary~\ref{coro: general index decom} gives a trivial zero-trace bilinear kernel and $\mathfrak b_1\equiv0$ on the two-dimensional compatible trace plane. Hence each of the cosine and sine sectors contributes two weak-kernel dimensions. For $n\ge2$, both the zero-trace kernel and the boundary kernel are trivial. Proposition~\ref{prop: second assumption steklov} therefore gives
\[
\operatorname{nul}_{\mathcal D}(Y_{\pi/6})
=\underbrace{1}_{n=0}
+\underbrace{2}_{\cos\theta}
+\underbrace{2}_{\sin\theta}
=5.
\]
\end{proof}

\end{document}
